\section[Cap\'{i}tulo III: \'{A}rea de Aplicaci\'{o}n]{\centering Cap\'{i}tulo III: \'{A}rea de Aplicaci\'{o}n}

\subsection{Introducci\'{o}n}

El presente cap\'{i}tulo describe el \'{a}rea donde se aplicar\'{a} el videojuego educativo de simulaci\'{o}n gamificada. En este documento no se desarrolla nuevamente el marco te\'{o}rico, sino que se muestra c\'{o}mo los conceptos de ciberseguridad, aprendizaje pr\'{a}ctico y gamificaci\'{o}n se concretan en un entorno de aplicaci\'{o}n para estudiantes de Inform\'{a}tica y Sistemas.

El proyecto se aplicar\'{a} como una herramienta complementaria para practicar identificaci\'{o}n y respuesta ante incidentes b\'{a}sicos de ciberseguridad. El prototipo no trabajar\'{a} sobre redes reales ni ejecutar\'{a} ataques reales; funcionar\'{a} mediante escenarios simulados donde el estudiante analiza pistas, toma decisiones y recibe retroalimentaci\'{o}n al finalizar cada caso.

\subsection{Contexto del \'{a}rea de aplicaci\'{o}n}

El \'{a}rea de aplicaci\'{o}n corresponde al aprendizaje pr\'{a}ctico de ciberseguridad en estudiantes de Ingenier\'{i}a Inform\'{a}tica e Ingenier\'{i}a de Sistemas de la Universidad Mayor de San Sim\'{o}n. En este contexto, la ciberseguridad requiere que el estudiante relacione conocimientos, habilidades y tareas, tal como plantea el marco NICE del NIST.

La necesidad de aplicaci\'{o}n surge porque las amenazas digitales no siempre se resuelven memorizando conceptos. El estudiante debe observar se\~{n}ales, comparar datos, reconocer comportamientos sospechosos y elegir respuestas adecuadas. Por ello, el videojuego se ubica como un recurso de pr\'{a}ctica guiada, seguro y repetible.

\begin{center}

\vspace{0.2cm}
\begin{tabular}{|p{0.30\textwidth}|p{0.55\textwidth}|}
\hline
\textbf{Elemento} & \textbf{Descripci\'{o}n} \\
\hline
Instituci\'{o}n & Universidad Mayor de San Sim\'{o}n. \\
\hline
Poblaci\'{o}n objetivo & Estudiantes de Ingenier\'{i}a Inform\'{a}tica e Ingenier\'{i}a de Sistemas. \\
\hline
Campo de aplicaci\'{o}n & Aprendizaje pr\'{a}ctico de ciberseguridad. \\
\hline
Tipo de herramienta & Videojuego educativo de simulaci\'{o}n gamificada. \\
\hline
Modalidad de uso & Prototipo monojugador con escenarios controlados. \\
\hline
\end{tabular}

\textbf{Tabla 3.1} - \textit{Contexto de aplicaci\'{o}n del prototipo}

\end{center}

\subsection{Usuarios e interesados}

Los usuarios principales del prototipo son los estudiantes, quienes interact\'{u}an directamente con los escenarios y toman decisiones frente a eventos sospechosos. Los docentes se consideran interesados secundarios, ya que pueden utilizar el videojuego como apoyo para reforzar contenidos y analizar resultados generales de desempe\~{n}o.

\begin{center}

\vspace{0.2cm}
\begin{tabular}{|p{0.25\textwidth}|p{0.32\textwidth}|p{0.28\textwidth}|}
\hline
\textbf{Usuario o interesado} & \textbf{Descripci\'{o}n} & \textbf{Responsabilidad o inter\'{e}s} \\
\hline
Estudiante & Usuario principal del videojuego. & Analizar eventos, tomar decisiones y revisar la retroalimentaci\'{o}n. \\
\hline
Docente & Usuario indirecto o de apoyo. & Usar el prototipo como recurso complementario para reforzar contenidos. \\
\hline
Instituci\'{o}n & Entorno acad\'{e}mico donde se ubica el proyecto. & Fortalecer recursos de formaci\'{o}n pr\'{a}ctica en ciberseguridad. \\
\hline
\end{tabular}

\textbf{Tabla 3.2} - \textit{Identificaci\'{o}n de usuarios e interesados}


\end{center}

\subsection{Descripci\'{o}n de la soluci\'{o}n propuesta}

La soluci\'{o}n propuesta consiste en un videojuego educativo donde el jugador asume el rol de analista de seguridad en un centro de operaciones simulado. Durante cada escenario, el jugador revisa correos, alertas, registros y comunicaciones sospechosas para decidir si corresponde aprobar, bloquear, reportar, escalar o contener un evento.

El escenario base se plantea en una organizaci\'{o}n ficticia, por ejemplo NovaTech S.A., con el objetivo de representar un contexto laboral sin comprometer infraestructura real. Las amenazas iniciales son \textit{phishing}, malware e ingenier\'{i}a social, debido a que pueden presentarse mediante pistas visibles y decisiones comprensibles para el estudiante. La relevancia del factor humano en incidentes de seguridad justifica este enfoque de simulaci\'{o}n.

\begin{center}

\vspace{0.2cm}
\begin{tabular}{|p{0.28\textwidth}|p{0.28\textwidth}|p{0.29\textwidth}|}
\hline
\textbf{Necesidad} & \textbf{Soluci\'{o}n propuesta} & \textbf{Resultado esperado} \\
\hline
Practicar identificaci\'{o}n de amenazas. & Escenarios con correos, alertas y llamadas sospechosas. & Mayor capacidad para reconocer indicios de riesgo. \\
\hline
Tomar decisiones ante incidentes. & Acciones de reportar, bloquear, escalar o contener. & Entrenamiento b\'{a}sico de respuesta ante eventos. \\
\hline
Recibir retroalimentaci\'{o}n clara. & Reporte de turno al finalizar cada escenario. & Comprensi\'{o}n de consecuencias de las decisiones. \\
\hline
Mantener un entorno seguro. & Simulaci\'{o}n sin redes, ataques ni sistemas reales. & Pr\'{a}ctica sin comprometer infraestructura. \\
\hline
\end{tabular}

\textbf{Tabla 3.3} - \textit{Soluci\'{o}n propuesta seg\'{u}n necesidad identificada}

\end{center}

\subsection{M\'{o}dulos del prototipo}

El prototipo se organiza en m\'{o}dulos que representan actividades frecuentes dentro de un escenario de seguridad. Esta organizaci\'{o}n permite separar las tareas del jugador y presentar informaci\'{o}n de forma progresiva.

\begin{center}

\vspace{0.2cm}
\begin{tabular}{|p{0.25\textwidth}|p{0.36\textwidth}|p{0.24\textwidth}|}
\hline
\textbf{M\'{o}dulo} & \textbf{Funci\'{o}n} & \textbf{Amenaza relacionada} \\
\hline
Mapa de casos & Organiza niveles y progresi\'{o}n de dificultad. & Todas las amenazas. \\
\hline
Bandeja de correo & Presenta mensajes con dominios alterados, urgencia falsa o adjuntos sospechosos. & \textit{Phishing} y malware. \\
\hline
Consola de eventos & Muestra registros, alertas y actividad an\'{o}mala. & Malware y respuesta a incidentes. \\
\hline
L\'{i}nea telef\'{o}nica & Representa solicitudes sospechosas de datos, c\'{o}digos o accesos. & Ingenier\'{i}a social. \\
\hline
Reporte de turno & Resume aciertos, errores y consecuencias. & Retroalimentaci\'{o}n del escenario. \\
\hline
\end{tabular}

\textbf{Tabla 3.4} - \textit{M\'{o}dulos principales del prototipo}

\end{center}

\subsection{Flujo de aplicaci\'{o}n}

El flujo de uso del prototipo inicia con la selecci\'{o}n de un caso. Luego, el jugador ingresa al escenario activo y revisa los eventos disponibles. Cada evento exige observar pistas antes de seleccionar una respuesta. Al finalizar, el sistema presenta un reporte con consecuencias.

\begin{center}

\vspace{0.2cm}
\begin{tabular}{|p{0.15\textwidth}|p{0.45\textwidth}|p{0.25\textwidth}|}
\hline
\textbf{Paso} & \textbf{Actividad del jugador} & \textbf{Resultado} \\
\hline
1 & Ingresar al prototipo y revisar el mapa de casos. & Selecci\'{o}n del escenario. \\
\hline
2 & Revisar correos, alertas o llamadas dentro del turno. & Identificaci\'{o}n de pistas. \\
\hline
3 & Elegir una acci\'{o}n: aprobar, bloquear, reportar, escalar o contener. & Registro de la decisi\'{o}n. \\
\hline
4 & Finalizar el caso. & Generaci\'{o}n del reporte de turno. \\
\hline
5 & Revisar consecuencias. & Retroalimentaci\'{o}n para mejorar decisiones futuras. \\
\hline
\end{tabular}

\textbf{Tabla 3.5} - \textit{Flujo general de uso}

\end{center}

\subsection{Relaci\'{o}n con el aprendizaje pr\'{a}ctico}

El \'{a}rea de aplicaci\'{o}n se vincula con el aprendizaje pr\'{a}ctico porque convierte conceptos de ciberseguridad en situaciones jugables. La gamificaci\'{o}n aporta retos, progresi\'{o}n y motivaci\'{o}n \parencite{kapp2012gamification}, mientras que la retroalimentaci\'{o}n permite que el estudiante comprenda la distancia entre su decisi\'{o}n y el desempe\~{n}o esperado.

Adem\'{a}s, las acciones propuestas se relacionan con fases b\'{a}sicas de manejo de incidentes, como identificaci\'{o}n, contenci\'{o}n y comunicaci\'{o}n \parencite{cichonski2022incident}. De esta manera, el videojuego no se limita a mostrar contenidos, sino que permite practicar una secuencia de razonamiento: observar, analizar, decidir y revisar consecuencias.

\subsection{S\'{i}ntesis del cap\'{i}tulo}

El \'{a}rea de aplicaci\'{o}n del proyecto se centra en el aprendizaje pr\'{a}ctico de ciberseguridad mediante un videojuego de simulaci\'{o}n. La propuesta se dirige a estudiantes de Inform\'{a}tica y Sistemas, utiliza un entorno ficticio tipo centro de operaciones de seguridad y organiza la experiencia mediante m\'{o}dulos, casos y reportes de retroalimentaci\'{o}n.

La estructura propuesta permite aplicar conceptos de \textit{phishing}, malware, ingenier\'{i}a social y respuesta a incidentes en un entorno seguro. Por tanto, el Cap\'{i}tulo III establece el espacio concreto donde el proyecto puede desarrollarse y validarse como herramienta educativa complementaria.
